\section{Normalized chord equations and passage through cones}
\label{sec:cones}

We work in the nonisotropic alternative of
\Cref{thm:isotropic-classification}. Thus the shape curve has a holomorphic
normal on its shape strip, the impact lift is the function
\(\mathfrak s\) of \Cref{thm:shape-lift}, and
\[
 f=\phi',\qquad q=f''/f,\qquad
 M(s)=\max\{|f(s)|,|q(s)|^{1/2}\}.
\]
Both \(f\) and \(q\) are holomorphic by \Cref{lem:f-zeros}.
The phase normalization is \eqref{shape:mean-normalization}:
\(\langle\mathfrak s-y\rangle=0\). In particular,
\(h_X(t)\le H\) by \eqref{shape:width-bound}. We first write the billiard equation in a form that
remains analytic when the affine coordinates of the shape curve have
poles. We then use its small-field estimate to continue the lift toward a
finite boundary of the shape strip.

\subsection{Kernel division and the exact recurrence}

Fix a shape point \(s\), a nonzero small parameter \(t\), and a normalized
step \(J\); the second endpoint is \(s+tJ\). The variable \(u\) below is a
fast coordinate, and derivatives with respect to \(u\) are denoted by
primes. Set
\begin{equation}\label{cone:fast-fields}
 F(u)=t f(s+tu),\qquad Q(u)=t^2q(s+tu),\qquad F''=QF,
\end{equation}
and introduce the three potentials
\begin{equation}\label{cone:kernel-potentials}
 Q_0=Q,\qquad Q_\pm=Q\mp\frac{3\ii}{2}F'-\frac{F^2}{4}.
\end{equation}
For each potential, let \(k_j(u)\) be the solution of
\(k_j''=Q_jk_j\) with \(k_j(0)=0\), \(k_j'(0)=1\).
The normalized arc-minus-chord action is
\begin{equation}\label{cone:kernel-action}
 p(s,J;t)=\frac{\text{arc minus chord}}{ct^3}
 =\frac{k_0(J)-\sqrt{k_+(J)k_-(J)}}{F(0)F(J)}.
\end{equation}
Here and below the square root is the branch continued from short real
chords at glancing. The letter \(p\) in this Section denotes this action,
not the offset of an unoriented line.

To verify \eqref{cone:kernel-action}, put
\[
 g_0=F,\qquad
 g_\pm(u)=F(u)\exp\left(\mp\frac{\ii}{2}\int_0^u F(v)\,\dd v\right).
\]
Direct differentiation, using \(F''=QF\), gives
\(g_j''=Q_jg_j\). Where \(g_j\ne0\), reduction of order gives
\begin{equation}\label{cone:reduction-order}
 k_j(J)=g_j(0)g_j(J)\int_0^Jg_j(u)^{-2}\,\dd u.
\end{equation}
The identity extends meromorphically across zeros. In the formula
\(\Gamma'=ce_\phi/f^2\), the arc integral is \(c\int f^{-2}\), and the
squared chord length is the product of the integrals of
\(cf^{-2}e^{\ii\phi}\) and \(cf^{-2}e^{-\ii\phi}\).
Substitution of \eqref{cone:reduction-order}, including the change
\(\dd s=t\,\dd u\), gives \eqref{cone:kernel-action} first for regular real
short chords and then by analytic continuation.

\begin{lemma}[Analytic kernel division]\label{lem:kernel-division}
Let \(Q\) be holomorphic on a fixed neighborhood of an integration
segment from \(0\) to \(J\). Regard \(Q\), \(J\), and the initial data
\(F(0),F'(0)\) for \(F''=QF\) as variables. Near zero initial data and a
pair \((Q,J)\) for which \(k_0(J)\ne0\), the right side of
\eqref{cone:kernel-action}, with the square root equal to \(k_0(J)\) at
zero initial data, extends holomorphically through
\(F(0)F(J)=0\). This extension is even in \(F\). Its bounds and analytic
dependence, including derivatives on smaller neighborhoods, are uniform
on compact sets of admissible \((Q,J)\).

In particular, for \(J\) in a fixed sufficiently small neighborhood of
\(1\), and small \(Q,F(0),F'(0)\),
\begin{equation}\label{cone:small-action}
 p-\frac{J^3}{24}
 =O\left(\|Q\|+|F(0)|^2+|F'(0)|^2\right).
\end{equation}
The same estimate holds for any fixed number of \(J\)-derivatives on a
smaller neighborhood.
\end{lemma}

\begin{proof}
The integral equation
\[
 k_j(u)=u+\int_0^u(u-v)Q_j(v)k_j(v)\,\dd v
\]
proves analytic dependence and uniform bounds by its uniformly convergent
Picard series. The same argument applies to \(F\) and to a moving
endpoint \(J\), provided the segments and their small thickenings remain
in the fixed domain.

Suppose first that \(F(0)=0\). If \(F\not\equiv0\), uniqueness gives
\(g_j=g_j'(0)k_j\), and
\(g_+'(0)g_-'(0)=F'(0)^2\). Since
\(g_+(J)g_-(J)=F(J)^2=F'(0)^2k_0(J)^2\), it follows that
\(k_+(J)k_-(J)=k_0(J)^2\). The conclusion also holds when \(F\equiv0\).
If instead \(F(J)=0\), apply the same argument to kernels started at
\(J\). The value at \(0\) of such a backward kernel is \(-k_j(J)\):
this follows by evaluating the constant Wronskian of the forward and
backward kernels at both endpoints. The two signs cancel in their
product, and the same identity follows.

At zero initial data, \(k_+k_-=k_0^2\ne0\). Consequently the square root
agreeing with \(k_0\) is analytic nearby. On each of the two zero
hyperplanes just considered it agrees with \(k_0\), by uniqueness of
that local branch. The map from initial data to endpoint data is
\[
 (F(0),F'(0))\longmapsto(F(0),F(J)),
\]
whose determinant is \(k_0(J)\). Its coordinates are therefore
independent. An analytic function vanishing on both coordinate
hyperplanes is divisible by their product, as is seen by its convergent
power series. This proves the claimed extension. Cauchy estimates in
these two coordinates prove locally uniform bounds for the quotient
and its derivatives, also when \(Q\) and \(J\) vary.

Replacing \(F\) by \(-F\) interchanges \(k_+\) and \(k_-\), so the
quotient is even. At \(F=Q=0\), its value is \(J^3/24\), as follows by
taking \(F\) to be a small constant in the formula computed below.
Evenness removes the terms linear in the initial data; analytic
interpolation from \(0\) to \(Q\) bounds the remaining constant term by
\(O(\|Q\|)\). This proves \eqref{cone:small-action} and its derivative
versions.
\end{proof}

For consecutive endpoints \(s_-\), \(s=s_-+tJ_-\), and \(s+tJ\),
stationarity of the sum of the two actions is exactly
\begin{equation}\label{cone:normalized-recurrence}
 p_J(s,J;t)-t p_s(s,J;t)=p_J(s_-,J_-;t).
\end{equation}
Indeed the derivative in the final endpoint of
\(ct^3p(s_-,(s-s_-)/t;t)\) is \(ct^2p_J\), whereas the derivative in the
initial endpoint of \(ct^3p(s,(s_+-s)/t;t)\) is
\(ct^3p_s-ct^2p_J\). Thus \eqref{cone:normalized-recurrence} has no
uncontrolled remainder. By \Cref{lem:kernel-division} it also makes
sense at the exceptional affine points of the curve.

\subsection{The two frozen actions}

The limits needed later are constant fast fields satisfying \(FQ=0\).
For \(F=\lambda\), \(Q=0\), set \(z=\lambda J/2\). Here
\(k_0(J)=J\) and \(k_\pm(J)=2\sin z/\lambda\). Consequently the action
and its energy \(E=Jp_J-p\) are
\begin{equation}\label{eq:frozen-F-energy}
 p_F(\lambda,J)=\frac{J}{\lambda^2}-\frac{2\sin z}{\lambda^3},
 \qquad
 E_F(\lambda,J)=\frac{2(\sin z-z\cos z)}{\lambda^3}.
\end{equation}
For \(F=0\), \(Q=\lambda^2\), the analytically divided action is
\begin{equation}\label{eq:frozen-Q-energy}
 p_Q(\lambda,J)=\frac{J}{2\lambda^2}-\frac{\tanh z}{\lambda^3},
 \qquad
 E_Q(\lambda,J)=\frac{\tanh z-z\operatorname{sech}^2z}{\lambda^3}.
\end{equation}
All quotients have their continuous values at \(\lambda=0\).
For the second calculation, take \(F=\eta e^{\lambda u}\), which solves
\(F''=\lambda^2F\), and expand the arc and chord in \(\eta\). The common
term of order \(\eta^{-2}\) cancels. More explicitly, put
\[
 A(u)=\int_0^u e^{\lambda v}\,\dd v,
 \qquad I_j=\int_0^J e^{-2\lambda u}A(u)^j\,\dd u.
\]
The two chord factors are
\(\eta^{-2}(I_0\pm\ii\eta I_1-\eta^2I_2/2+O(\eta^3))\).
Their product is
\(\eta^{-4}(I_0^2+\eta^2(I_1^2-I_0I_2)+O(\eta^4))\), so subtraction
from the arc \(\eta^{-2}I_0\), with the glancing square root, leaves
\((I_2-I_1^2/I_0)/2+O(\eta^2)\). Since
\(A(u)=(e^{\lambda u}-1)/\lambda\), the limit is
\[
 \frac1{2\lambda^2}
 \left(J-\frac{\left(\int_0^J e^{-\lambda u}\,\dd u\right)^2}
 {\int_0^J e^{-2\lambda u}\,\dd u}\right)
 =\frac{J}{2\lambda^2}-\frac{\tanh(\lambda J/2)}{\lambda^3}.
\]
Lemma~\ref{lem:kernel-division} identifies this limit with the analytic
extension at \(F=0\).

For later use, the expansions at zero are
\begin{align}
 p_F&=\frac{J^3}{24}-\frac{\lambda^2J^5}{1920}+O(\lambda^4J^7),&
 E_F&=\frac{J^3}{12}-\frac{\lambda^2J^5}{480}+O(\lambda^4J^7),
 \label{cone:F-series}\\
 p_Q&=\frac{J^3}{24}-\frac{\lambda^2J^5}{240}+O(\lambda^4J^7),&
 E_Q&=\frac{J^3}{12}-\frac{\lambda^2J^5}{60}+O(\lambda^4J^7).
 \label{cone:Q-series}
\end{align}

\begin{lemma}[Convergence to a frozen action]\label{cone:profile-action-convergence}
Fix a compact path of frozen states from the small-field germ to a
finite point of one of the two displayed actions above, avoiding kernel
and root degeneracies, with \(J\ne0\) and \(p_{JJ}\ne0\) at its endpoint.
Suppose analytic fast fields satisfying \(F''=QF\) converge locally
uniformly to those frozen states on complex neighborhoods of the path
and of all its swept integration segments. Continue their glancing
square-root germs along
the corresponding perturbed path. Their actions then converge
analytically, with all derivatives on smaller neighborhoods of the
endpoint, to the frozen action on that fixed continued sheet. The
endpoint alone does not select a sheet after an arbitrary loop.
\end{lemma}

\begin{proof}
In the first case,
\(p_{JJ}=\sin z/(2\lambda)\); hence its nonvanishing excludes the zero of
\(k_\pm\), while \(k_0=J\ne0\). The ordinary quotient formula and
analytic dependence of kernels apply. In the second case,
\[
 p_{JJ}=\frac{\tanh z\,\operatorname{sech}^2z}{2\lambda},\qquad
 k_0=k_+=k_-=\frac{\sinh(2z)}{\lambda}.
\]
A zero of this last kernel either has \(\sinh z=0\), when
\(p_{JJ}=0\), or \(\cosh z=0\), when the displayed branch is not finite.
Thus the division lemma applies. The limiting case \(\lambda=0\) has
\(p_{JJ}=J/4\). A compact path can be covered by finitely many of these
analytic neighborhoods. Uniqueness of the square root on each overlap
tracks the chosen branch, and Cauchy estimates give convergence of all
derivatives.
\end{proof}

\subsection{Continuation in cones of arbitrary fixed aperture}

For this subsection suppose \(H<\infty\). Depth always means distance
below the upper edge:
\[
 \delta=H-\Im y,\qquad \delta_s=H-\Im s.
\]
The lower edge is treated by real symmetry.

\begin{theorem}[Cone passage]\label{thm:cone}
Fix \(1/2<\alpha<1\), a finite target aperture \(m_0\), and a fixed
\(\delta_b\in(0,H/2)\). There is an input aperture \(L\), depending on
\(\alpha\) and \(m_0\), with the following property. Fix \(K<\infty\),
let \(d\to0\), choose real centers \(x_0\), and let \(0<r\le Kd\).
Suppose
\begin{equation}\label{cone:input}
 rM(s)\le\eps(d/\delta_s)^\alpha
 \quad\text{if}\quad
 d/10\le\delta_s\le\delta_b,
 \quad |\Re s-x_0|\le L\delta_s.
\end{equation}
For sufficiently small \(\eps>0\) and then sufficiently small \(d\),
the lift continues jointly holomorphically along vertical paths from a
fixed lower row throughout
\begin{equation}\label{cone:target}
 d<\delta<\delta_*,\qquad
 |\Re y-x_0|<m_0\delta,\qquad
 |t|<r(\delta/d)^\alpha,
\end{equation}
where \(\delta_*>0\) is fixed. With
\(B=1+(\eps/r)d^\alpha\), it satisfies
\begin{equation}\label{cone:passage-estimate}
 |\mathfrak s(y,t)-y|
 \le C B^2|t|^2\delta^{1-2\alpha}.
\end{equation}
The relative depth error is small, so these continuations remain in the
shape strip, and they agree with the original lift wherever the latter
is defined below them.

The input aperture can be chosen independently of \(K\). Once it
contains the fixed geometric enlargements described in the proof, the
constant \(C\) and the smallness requirement on \(\eps\) can be chosen
independently of the target aperture. They may depend on \(K\),
\(\alpha\), and the fixed lower-strip data.
\end{theorem}

\begin{figure}[t]
\centering
\begin{tikzpicture}[>=Stealth,x=1.3cm,y=1.1cm]
\fill[deepblue!8] (-.12,2.85)--(-2.45,0)--(2.45,0)--(.12,2.85)--cycle;
\draw[thick] (-3,3)--(3,3) node[right,font=\small] {$\Im s=H$};
\draw[deepblue,dashed] (-2.45,0)--(0,3)--(2.45,0);
\draw[gray] (-2.6,.30)--(2.6,.30);
\node[font=\small,anchor=west] at (2.65,.30) {starting slab};
\draw[gray,dotted] (-2.0,2)--(2.0,2);
\fill (0,2) circle (1.7pt);
\node[font=\small,anchor=west] at (2.10,2) {$x+\ii(H-d)$};
\draw[<->] (1.62,2.04)--node[right,font=\small] {$d$} (1.62,2.96);
\draw[->,warmred,thick] (-.65,.35)--(-.49,.75);
\draw[->,warmred,thick] (-.49,.75)--(-.33,1.15);
\draw[->,warmred,thick] (-.33,1.15)--(-.17,1.55);
\draw[->,warmred,thick] (-.17,1.55)--(0,2);
\node[font=\small,anchor=east,warmred] at (-2,1.12) {steps $\pm t$};
\node[font=\small,align=center] at (0,-.42)
 {$|\Re s-x|\le L(H-\Im s)$\qquad with a fixed enlarged margin};
\end{tikzpicture}
\caption{A cone passage begins in known analytic data and follows steep
increments toward the chosen depth. Local error bounds are summed in
depth, with exponent greater than \(1/2\). Enlarging a fixed aperture
provides geometric margin without changing the local constants.
Parameter disks subsequently fill the directions close to real increments.}
\label{fig:cone}
\end{figure}


\begin{proof}
We give separately the recurrence argument in steep parameter
directions, illustrated in \Cref{fig:cone}, and the analytic-disk argument that fills the remaining
directions.

Fix a small number \(\gamma>0\). Suppose first that
\(|\Im t|\ge\gamma|t|>0\). By evenness in \(t\), use whichever of
\(t,-t\) moves upward. Start in a fixed slab about depth
\(\delta_b/2\), where joint analyticity and the phase normalization give
\[
 \mathfrak s-y=O(t^2),\qquad
 \frac{\mathfrak s(y+t,t)-\mathfrak s(y,t)}{t}=1+O(t^2)
\]
with derivative bounds uniform in the real coordinate modulo its
period. The largest parameters in the construction tend to zero,
since \(r d^{-\alpha}\le Kd^{1-\alpha}\to0\).

Follow the exact recurrence toward the target, denoting successive
unperturbed depths by \(x_j\). They decrease by \(|\Im t|\). We allow
here the slightly larger range
\(x_j\ge d/2\), \(|t|\le10r(x_j/d)^\alpha\), which will be useful in
the disk construction. Suppose inductively that position errors are a
small fraction of the depths and that normalized steps are close to
\(1\). The shape fields on every required segment and a small
thickening satisfy
\[
 |F(0)|+|F'(0)|\le C_K|t|Bx_j^{-\alpha},\qquad
 \|Q\|\le C_K|t|^2B^2x_j^{-2\alpha}.
\]
To see the derivative estimate, apply Cauchy on a disk of radius a
small fixed fraction of \(x_j\). Although \(|t|/x_j\) need not tend to
zero, it is bounded by a constant depending on \(K\). Consecutive
depths are therefore comparable, with constants depending on \(K\),
and the thickenings may be taken with smaller such fractions. They
remain above depth \(d/10\). Their horizontal displacements lie in the
direction cone of slope at most \(1/\gamma\), with a fixed extra
margin. Increasing the input aperture by this fixed geometric amount
suffices; reducing the thickness when \(K\) increases does not require
a larger aperture.

The small parameter in the kernel estimates is bounded by
\[
 |t|Bx_j^{-\alpha}
 \le C\bigl(\eps+rd^{-\alpha}\bigr).
\]
Thus \Cref{lem:kernel-division} applies uniformly. If
\(a(s,J;t)=p(s,J;t)-J^3/24\), then
\begin{equation}\label{cone:action-estimates}
 |\partial_J^\nu\partial_s^\mu a|
 \le C_{K,\mu,\nu}|t|^2B^2x_j^{-2\alpha-\mu}
\end{equation}
for the finitely many derivatives used below. The \(s\)-derivative
bounds follow by Cauchy on the same thickenings.

Evaluate \eqref{cone:normalized-recurrence} with the trial new step
\(J=J_-\). The cubic contributions cancel. The difference of
\(a_J\) at the two adjacent basepoints, together with \(ta_s\), is at
most
\(C_K|t|^3B^2x_j^{-2\alpha-1}\). The derivative of the equation in
the new \(J\) is \(J/4+o(1)\), uniformly bounded away from zero.
The analytic implicit function theorem therefore supplies a unique
nearby step, with
\begin{equation}\label{cone:one-step}
 |J-J_-|\le C_K|t|^3B^2x_j^{-2\alpha-1}.
\end{equation}
For \(\beta>1\), comparison with an integral on the depth mesh gives
\[
 |t|\sum_{x_j\ge x}x_j^{-\beta}\le C_{K,\gamma,\beta}x^{1-\beta}.
\]
The possible final summand is controlled because \(|t|/x\) is bounded.
Summing \eqref{cone:one-step} first, and then summing the position
increments \(t(J-1)\), gives
\begin{equation}\label{cone:steep-estimates}
 |J-1|\le C_K|t|^2B^2x^{-2\alpha},\qquad
 |\mathfrak s-y|\le C_K|t|^2B^2x^{1-2\alpha}.
\end{equation}
The second sum uses \(2\alpha>1\). Initial errors from the fixed slab
are absorbed because \(B\ge1\). Moreover
\[
 \frac{|\mathfrak s-y|}{x}+|J-1|
 \le C_K\bigl(\eps+rd^{-\alpha}\bigr)^2,
\]
which closes the induction after the prescribed order of choices.

These recurrences define analytic germs in \((y,t)\). Locally choose a
fixed integer number of steps back into the starting slab, where both
initial neighbors are known. Two such choices agree by uniqueness of
the near-identity recurrence. They consequently give analytic
continuation from below on every steep sector under discussion.

We now fill all parameter directions. Call a germ reached if it is
obtained by continuing from the fixed lower row along the vertical path
with \(\Re y\) and \(t\) fixed. Reached germs form a single-valued
holomorphic function on their open domain: nearby parallel paths are
compared through a finite chain of the same germs from below. The steep
continuations just constructed agree with these germs. Work with
\[
 U(y,t)=\frac{\mathfrak s(y,t)-y}{t^2},
\]
which is regular at \(t=0\) on the initial domain and on its
continuations.

Choose a large fixed \(D_0\) with \(D_0^\alpha>10\), then a sufficiently
small fixed \(\gamma\). For a target \((y_c,t_c)\) of depth \(\delta\),
put \(R=r(\delta/d)^\alpha\). Since \(|t_c|<R\), the circle
\(t=t_c+3R\zeta\), \(|\zeta|=1\), meets the nonsteep sector only on
two uniformly short arcs. There is a polynomial \(g\), with
\(g(0)=0\), such that
\begin{equation}\label{cone:disk-bending}
 |g(\zeta)|\le C_0\delta,
 \quad \Im g\le\delta/4\ \text{on the circle},
 \quad \Im g\le-D_0\delta\ \text{on those arcs and a collar}.
\end{equation}
Here \(C_0\) is fixed independently of the target aperture. For
completeness, choose smooth negative bumps equal to a sufficiently
large negative constant on the two arcs and supported in slightly
larger short intervals. Add a constant to make their mean zero.
The intervals can be made so short that the added constant is less
than \(1/8\). Approximate the resulting smooth function, with a fixed
strict margin in its inequalities, by a real trigonometric polynomial
of mean zero. It is the boundary imaginary part of an analytic
polynomial vanishing at zero. Multiplication by \(\delta\) gives
\eqref{cone:disk-bending}. The centers of the arcs and
\(t_c/(3R)\) range over a compact set with modulus at most \(1/3\),
so the bump widths and polynomial degree can be fixed uniformly.
This also bounds the conjugate real part by \(C_0\delta\).

Increase an auxiliary target aperture \(m\ge m_0\) so that
\[
 m+C_0<mD_0.
\]
Choose \(L\) to contain this aperture, the shifts by \(C_0\delta\),
and the steep paths and thickenings above, all with fixed slack. We
induct from a sufficiently small fixed depth down through successive
bands whose depth thresholds are halved. In a new band use the disks
\begin{equation}\label{cone:disks}
 t=t_c+3R\zeta,\qquad y=y_c+g(\zeta)+w,
 \qquad |\zeta|\le1,
\end{equation}
where \(w\) is a small complex parameter. On the nonsteep boundary
arcs the depth is at least \(D_0\delta\). They are therefore in the
previous all-direction domain; the inequality
\(D_0^\alpha>10\) supplies the parameter margin, and
\(m+C_0<mD_0\) supplies the aperture margin. On the remaining arcs,
steep continuation applies at depth at least \(\delta/2\), with the
larger parameter range already allowed. Slight collars of the arcs
have the same properties.

Start \(y_c\) lower by a fixed sufficiently large multiple of
\(\delta\); then the entire disk and a neighborhood lie in the
previous domain. Raise \(y_c\) to its target value in small steps.
The holomorphic disk-continuation principle in
\Cref{lem:meromorphic-hartogs,cont:disc-sweep}, with the parameter \(w\),
continues \(U\) over the disks and identifies the resulting central
germ with vertical continuation from below.

There is no growth of the estimate with the number of bands. If the
inductive bound is \(A B^2x^{1-2\alpha}\), its value on a previous-domain
arc of depth \(x\ge D_0\delta\) is at most
\(A D_0^{1-2\alpha}B^2\delta^{1-2\alpha}\), with the fixed factor
\(D_0^{1-2\alpha}<1\). On a steep arc
\eqref{cone:steep-estimates} gives an independent constant times
\(B^2\delta^{1-2\alpha}\). Choose \(A\) once larger than that constant
and the initial lower-strip constant. Maximum modulus on each disk
retains the same \(A\). Starting the induction at a sufficiently small
fixed depth keeps its lower dips below \(\delta_b\); a partial last
band reaches any \(\delta>d\). This proves
\eqref{cone:passage-estimate}.

The constants in the depth sums, field bounds, and maximum-modulus
argument depend only on the fixed local margins. They do not depend
on how far the cone extends horizontally. Initial data are uniform
modulo the real period. Thus increasing a target aperture only
increases the input aperture needed to contain the indicated
geometric enlargements; it does not enlarge the estimate constant or
change the smallness choice. The relative-error estimate already
proved keeps the image below the upper shape edge and, after reducing
\(\delta_*\), above the lower edge.
\end{proof}

\begin{remark}[Reference heights]\label{cone:reference-height}
The local construction in \Cref{thm:cone} does not use the assertion that
\(H\) is a singularity height. It remains valid with a reference height
\(h\) in place of \(H\), provided the stated cone and its thickenings
are contained in the actual shape domain, lower-row initial data are
available, and the relative-error bound keeps the constructed lift in
that domain. In a rescaled sequence this formulation is applied on
fixed compact sets, with thresholds allowed to depend on the chosen
finite aperture. The constants in the estimate remain independent of
that aperture. A statement about the global width still uses the
actual first boundary of the phase domain, as explained next.
\end{remark}

\subsection{Quiet profiles and their consequences}

At a real parameter \(t\), any passage in \Cref{thm:cone} gives a global
lower bound on \(h_X(t)\). Indeed, if the original phase strip ended at
an earlier height, the local identity
\(X(y,t)=\Gamma(\mathfrak s(y,t))\), valid jointly in \((y,t)\) and with
\(|\Im\mathfrak s|<H\), would give a joint meromorphic patch through
that first edge. This contradicts \Cref{thm:no-patch}. This argument
uses the no-patch assertion only at a real physical parameter.

\begin{corollary}[Coarse bounds and quiet profiles]\label{cor:quiet-cone}
Let \(d_n\to0\), \(b_n\ge1\), and \(x_n\in\R\). Suppose, on a
sufficiently large fixed cone starting at depth \(d_n\), that
\begin{equation}\label{cone:coarse-bound}
 \frac{d_n}{b_n}M(s)\le C(d_n/\delta_s)^\alpha,
 \qquad \frac12<\alpha<1.
\end{equation}
The following statements hold; whenever an arbitrary fixed aperture is
used, the hypothesis is required on its corresponding fixed
enlargement.
\begin{enumerate}[label=\textup{(\roman*)}]
\item On a strictly coarser cone, starting at a fixed sufficiently large
multiple of \(d_n\), there is a passage with parameter radius
\(c(d_n/b_n)(\delta/d_n)^\alpha\), and
\begin{equation}\label{cone:coarse-error}
 |\mathfrak s-y|
 =O\left(b_n^2|t|^2d_n^{2\alpha-2}\delta^{1-2\alpha}\right).
\end{equation}
Equivalently, for
\(Y=(y-x_n-\ii H)/d_n\),
\(S=(\mathfrak s-x_n-\ii H)/d_n\), and \(T=tb_n/d_n\),
\[
 |S-Y|\le C|T|^2(-\Im Y)^{1-2\alpha}
\]
far enough down at each bounded \(T\).
\item Suppose in addition that
\begin{equation}\label{cone:quiet-profile}
 \frac{d_n}{b_n}M(x_n+\ii H+d_n S)\longrightarrow0
\end{equation}
uniformly on every compact subset of \(\Im S<0\).
For every \(c_1,T_1>0\), lift passages reach depth \(c_1d_n\) for all
\(|t|\le T_1d_n/b_n\), for large \(n\). At real parameters, with
\(\Delta(t)=H-h_X(t)\),
\begin{equation}\label{cone:quiet-width}
 \sup_{|T|\le T_1}\frac{\Delta(Td_n/b_n)}{d_n}\longrightarrow0.
\end{equation}
\item The hypotheses in \textup{(ii)} are impossible if \(b_n\) stays
bounded. In fact a vanishing weighted cone bound down to a sufficiently
small fixed positive depth ratio already gives this contradiction.
\end{enumerate}
\end{corollary}

\begin{proof}
For \textup{(i)}, apply \Cref{thm:cone} with its starting scale a fixed
large multiple of \(d_n\) and with \(r\) a sufficiently small fixed
multiple of \(d_n/b_n\). Then its \(B\) is bounded by a constant times
\(b_nd_n^{\alpha-1}\), which gives \eqref{cone:coarse-error}. Increasing
the fixed starting multiple permits any prescribed bounded interval
of \(T\).

For \textup{(ii)}, choose \(\beta\) with \(1/2<\beta<\alpha\).
On a rescaled cone put \(D=\delta_s/d_n\). The weighted norm to be made
small is
\(D^\beta(d_n/b_n)M\). On a compact range of \(D\) it tends to zero
by \eqref{cone:quiet-profile}; on the tail \(D\ge R\),
\eqref{cone:coarse-bound} bounds it by \(CR^{\beta-\alpha}\).
First take \(R\) large, then \(n\) large. This proves a vanishing
weighted bound down to every fixed positive depth ratio. Apply
\Cref{thm:cone} with a small fixed multiple of \(c_1d_n\) as its scale
and a large enough fixed multiple of \(T_1d_n/b_n\) as its radius.
Its ratio \(r/d\) is bounded because \(b_n\ge1\). The resulting
relative error tends to zero. The real no-patch transfer gives
\(h_X(t)\ge H-c_1d_n\) uniformly in the stated real interval.
Since \(c_1\) is arbitrary and \(h_X\le H\), this proves
\eqref{cone:quiet-width}.

For \textup{(iii)}, boundedness of \(b_n\) turns the vanishing weighted
bound into the same statement with normalization \(d_nM\). Use
\Cref{thm:cone} with starting scale \(d_n/20\), radius \(2d_n\), and
target aperture at least \(50\). Its hypotheses need only hold down
to a fixed smaller depth ratio. We obtain the shifted impact data
needed on the boundaries of the disks
\begin{equation}\label{cone:quiet-disks}
 t=d_n\zeta,\qquad
 y=x_n+\ii\sigma-\frac{\ii d_n}{4}\zeta^2+w,
 \quad |\zeta|\le1,
\end{equation}
where \(w\) is small and \(\sigma\) rises from \(H-4d_n\) to \(H\).
On upper imaginary parameter arcs, recover \(L(y,t)\) from the
preceding two impacts, at phases \(y-t/2\) and \(y-3t/2\), and
reflection at the nearer impact. On lower arcs use the two forward
impacts. These are meromorphic identities in the line and point data,
so exceptional affine configurations are covered by continuation.
At \(\sigma=H\), \(w=0\), and \(\zeta=e^{\ii\theta}\), the nearer
of the two selected impacts has depth
\[
 d_n\left(\frac{1-2\sin^2\theta}{4}
                 +\frac{|\sin\theta|}{2}\right)\ge\frac{d_n}{4}.
\]
The farther one is deeper, and all horizontal shifts fit the fixed
aperture. On overlaps near real parameter directions both
constructions agree with the primary line by continuation from below.
At the starting height the full disks have the necessary impact data.
The meromorphic disk-continuation principle therefore fills
\eqref{cone:quiet-disks} up to \(\sigma=H\). At their centers it gives
a joint patch of \(L\) through \((x_n+\ii H,0)\), attached to the
primary zero-parameter data, contrary to \Cref{thm:no-patch}.
\end{proof}

\begin{proposition}[Subcritical growth is impossible]\label{prop:subcritical-exclusion}
At a finite shape edge, a uniform bound
\(\max_{\Im s=H-\delta}M(s)=O(\delta^{-p_0})\) with \(p_0<1\) is
impossible in the nonisotropic alternative.
\end{proposition}

\begin{proof}
Choose \(\alpha\) with \(\max\{p_0,1/2\}<\alpha<1\). On any cone
with \(\delta_s\ge d\),
\[
 dM(s)\le C d^{1-p_0}(d/\delta_s)^\alpha
                 (\delta_s/d)^{\alpha-p_0}.
\]
More directly, the weighted quantity
\(dM(s)(\delta_s/d)^\alpha\) is bounded by
\(C d^{1-\alpha}\delta_s^{\alpha-p_0}\), which tends to zero
uniformly for depths between any fixed multiple of \(d\) and the
fixed starting row. Thus \Cref{cor:quiet-cone}\textup{(iii)} applies
with \(b=1\), giving a contradiction.
\end{proof}

\begin{proposition}[A witness near every center]\label{cone:no-quiet-witness}
Suppose, for all sufficiently small \(\delta_s\), that
\[
 \delta_s\max_{\Im s=H-\delta_s}M(s)
 \le C\bigl(1+\log(1/\delta_s)\bigr)^A
\]
for some fixed \(C,A\). Then there are constants
\(L,C_1,C_2,c_*>0\) such that for every sufficiently large
\(k=\log(1/d)\) and every real center \(x_0\), there is a point \(s\)
with
\begin{equation}\label{cone:witness-window}
 |\Re s-x_0|\le L\delta_s,\qquad
 k-C_1\log k\le\log(1/\delta_s)\le k+C_2,
 \qquad \delta_sM(s)\ge c_*.
\end{equation}
\end{proposition}

\begin{proof}
Fix \(1/2<\alpha<1\) and the cone aperture needed in the bounded-scale
quiet contradiction. Choose \(C_2\) so that \(e^{-C_2}d\) is below
all the fixed depth ratios required there, and then choose
\(C_1>A/(1-\alpha)\). If no uniform \(c_*\) exists, there are
\(k_n\to\infty\) and centers such that
\(\sup\delta_sM(s)\to0\) in their windows. On that window,
\[
 d_nM(s)(\delta_s/d_n)^\alpha
 =\delta_sM(s)(d_n/\delta_s)^{1-\alpha}\longrightarrow0
\]
uniformly, because the last factor is at most
\(e^{(1-\alpha)C_2}\). Below the window, where
\(\delta_s/d_n\ge k_n^{C_1}\), the global polynomial bound gives
\[
 \delta_sM(s)(d_n/\delta_s)^{1-\alpha}
 \le C'(1+k_n)^A k_n^{-C_1(1-\alpha)}\longrightarrow0.
\]
The remaining fixed lower-strip portion has bounded \(M\) and is
smaller still. This is the vanishing cone bound forbidden by
\Cref{cor:quiet-cone}\textup{(iii)}, with \(b_n=1\).
\end{proof}

\subsection{Exclusion of finite superlinear growth}

\begin{theorem}[Exclusion of finite superlinear growth]\label{thm:finite-superlinear-exclusion}
Suppose \(H<\infty\), and define
\[
 g(k)=\log\max_{\Im s=H-e^{-k}}M(s).
\]
In the nonisotropic alternative it is impossible that
\begin{equation}\label{cone:finite-order-assumption}
 1<\limsup_{k\to\infty}\frac{g(k)}{k}<\infty.
\end{equation}
\end{theorem}

\begin{proof}
Write the limsup in \eqref{cone:finite-order-assumption} as \(\rho\),
and fix
\[
 1<\alpha_1<\alpha_2<\rho<P.
\]
We first select scales with useful bounds on both sides. There are
\(k_n\to\infty\) such that \(g(k_n)\ge\alpha_2k_n\) and
\begin{equation}\label{cone:two-sided-envelopes}
 g(k)\le g(k_n)+
 \begin{cases}
 \alpha_1(k-k_n),&k_0\le k\le k_n,\\
 P(k-k_n),&k\ge k_n.
 \end{cases}
\end{equation}
Here \(k_0\) is fixed and sufficiently large. To prove this, take
arbitrarily late running records \(r_n\) of \(g(k)-\alpha_1k\) for
which \(g(r_n)/r_n>\alpha_2\). Such records exist: a time with
\(g(k)>\alpha_2k\) forces the preceding running maximum to occur at a
time \(r\le k\) with
\(g(r)-\alpha_1r>(\alpha_2-\alpha_1)k\ge(\alpha_2-\alpha_1)r\).
The maximizing times tend to infinity. Since \(P>\rho\), the function
\(g(k)-Pk\) tends to \(-\infty\); let \(k_n\ge r_n\) maximize it on
\([r_n,\infty)\). For \(k\ge k_n\) this gives the second inequality
in \eqref{cone:two-sided-envelopes}. If \(r_n\le k\le k_n\), it gives
\(g(k)\le g(k_n)+P(k-k_n)\le g(k_n)+\alpha_1(k-k_n)\).
For \(k\le r_n\), combine the record bound with
\(g(k_n)\ge g(r_n)+P(k_n-r_n)\), obtaining the first inequality again.
The same last inequality gives \(g(k_n)\ge\alpha_2 k_n\).

Set
\begin{equation}\label{cone:superlinear-scales}
 d_n=e^{-k_n},\qquad m_n=e^{g(k_n)},\qquad b_n=d_nm_n.
\end{equation}
Then \(b_n\to\infty\), and in fact
\begin{equation}\label{cone:b-dominates-log}
 \frac{b_n}{\log m_n}\longrightarrow\infty,
\end{equation}
because \(b_n\ge e^{(\alpha_2-1)k_n}\), whereas
\(\log m_n=O(k_n)\).

\medskip\noindent
\emph{Compact profiles and a trivial pair.}
For any choice of real centers \(x_n\), introduce the rescaled shape
coordinate \(S=(s-x_n-\ii H)/d_n\) and fields
\begin{equation}\label{cone:superlinear-profiles}
 F_n(S)=\frac{f(x_n+\ii H+d_nS)}{m_n},\qquad
 Q_n(S)=\frac{q(x_n+\ii H+d_nS)}{m_n^2}.
\end{equation}
Their domains exhaust the lower half-plane. The envelopes give,
with \(D=-\Im S\),
\begin{equation}\label{cone:profile-envelope}
 \max\{|F_n(S)|,|Q_n(S)|^{1/2}\}
 \le
 \begin{cases}
 D^{-\alpha_1},&D\ge1,\\
 D^{-P},&0<D\le1,
 \end{cases}
\end{equation}
on every fixed compact set for all large \(n\). These bounds are
uniform in all real centers. Montel's theorem gives subsequential
holomorphic limits \(F,Q\). The rescaled differential equation is
\begin{equation}\label{cone:profile-ode}
 F_n''=b_n^2Q_nF_n.
\end{equation}
Cauchy estimates bound the left side locally, so \(FQ=0\). Since the
half-plane is connected, either \(F\equiv0\) or \(Q\equiv0\).
Every limit also satisfies the deeper bound \(D^{-\alpha_1}\) in
\eqref{cone:profile-envelope}.

We claim that some sequence of centers, after passage to a subsequence,
has \(F=Q=0\). Suppose otherwise. Fix a small closed disk about \(-\ii\)
inside the half-plane and take the supremum norms of the two fields on
it. Compactness then gives a constant \(c>0\) such that, for all centers
and large \(n\), the larger of the two norms is at least \(c\).
Otherwise a sequence with both norms tending to zero would have the
zero pair as its limit everywhere, by analytic uniqueness. Purity and
compactness likewise imply that the smaller norm tends to zero
uniformly over centers: any contrary sequence would have a limiting
pair with both norms positive. The norms depend continuously on the
center. Once their smaller value is less than \(c/3\), the two
possible dominant types give disjoint open sets covering the connected
circle of centers. Thus one type holds at every center. Passing to a
subsequence, it is the same type for every \(n\).

Suppose first that it is the \(F\)-type. All local limiting functions
\(F\) obtained from arbitrary centers are nonzero. This compact family
has a uniform bound on the number of zeros in any fixed smaller disk
about \(-\ii\), and, for every fixed small \(\eta>0\), a uniform
positive lower bound away from the \(\eta\)-neighborhoods of those
zeros. Here is a direct compactness justification. A convergent
subsequence has a nonzero holomorphic limit; choose a slightly larger
circle without zeros of that limit. The argument principle bounds the
number of zeros inside, and uniform convergence bounds the function
below on a smaller compact set away from the zeros. A failure of
either uniform assertion would produce a contrary convergent
subsequence. The same assertions hold on the slightly enlarged scaled
neighborhoods needed to cover a narrow band about the row \(D=1\).

Cover the physical periodic row at depth \(d_n\) by \(O(d_n^{-1})\)
cells of size comparable to \(d_n\). Around its nearby zeros remove
disks of radii between \(\eta d_n\) and \(2\eta d_n\), with generic
radii to avoid boundary tangencies. Choose \(\eta\) sufficiently small
in terms of the uniform local zero count. Every connected cluster
meeting the row then has a uniformly bounded number of disks and
diameter much smaller than \(d_n\). Indeed, a minimal chain escaping a
fixed enlarged cell would require more zeros in that cell than the
uniform bound allows. Fill holes in the clusters and discard nested
obstacles. Detour the row along the outer boundary arcs of these
clusters. The resulting closed contour \(C_n\) winds once around the
periodic cylinder, stays in a fixed narrow band about depth \(d_n\),
and satisfies
\begin{equation}\label{cone:large-f-contour}
 \operatorname{length}(C_n)=O(1),\qquad
 |f|\ge c_\eta m_n\quad\hbox{on }C_n.
\end{equation}
The length estimate follows because the total number of relevant
zero disks is \(O(d_n^{-1})\), each with circumference \(O(d_n)\).

All poles of \(f^{-2}\) in the shape strip are at simple zeros of
\(f\). At such a zero \(z\), \Cref{lem:f-zeros} gives
\(f''(z)=0\), and hence
\[
 \operatorname{Res}_{z}(f^{-2})
 =-\frac{f''(z)}{f'(z)^3}=0.
\]
Deforming \(C_n\) to the real period, with residues included, therefore
gives
\[
 \int_{C_n}\frac{\dd s}{f(s)^2}
 =\int_0^{2\pi}\frac{\dd s}{f(s)^2}>0.
\]
The right side is a fixed positive number because \(f\) is positive on
the real curve. The absolute value of the left side is
\(O(m_n^{-2})\) by \eqref{cone:large-f-contour}. This is a
contradiction.

Suppose next that the common type is the \(Q\)-type. There are centers
\(x_n\) with
\begin{equation}\label{cone:f-lower-center}
 |f(x_n+\ii(H-d_n))|\ge c_0>0.
\end{equation}
Indeed the horizontal mean of \(\log|f|\) is convex as a function of
height, by subharmonicity and periodicity. It is even by real symmetry,
so it is at least its finite real-row value. At least one point on
each row has modulus no smaller than the exponential of that mean.
Zeros cause no difficulty because \(\log|f|\) is locally integrable.
At these centers pass to a subsequence with \(Q_n\to Q\not\equiv0\).

Near a point where \(Q\ne0\), choose a unit complex number \(e\) and a
fixed short segment in direction \(e\) on which
\(\Re(e^2Q_n)\ge c'>0\) for all large \(n\). On that segment, with
real parameter \(v\), the nonnegative function
\(u(v)=|F_n(z+ev)|^2\) satisfies
\[
 u''=2|eF_n'|^2+2b_n^2\Re(e^2Q_n)|F_n|^2
 \ge 2c'b_n^2u.
\]
The endpoint values are uniformly bounded by normality. Comparison
with the solution of \(v''=2c'b_n^2v\) having those endpoint bounds
shows, on a shorter central segment,
\(u\le C e^{-c b_n}\). To justify comparison, subtract that solution;
a positive interior maximum of the difference contradicts
\(w''-2c'b_n^2w\ge0\). The same estimate holds uniformly on a small
neighborhood of the chosen point.

Choose a circle centered at \(-\ii\), contained in the half-plane,
and avoiding the discrete zero set of \(Q\). A finite cover by these
neighborhoods gives \(|F_n|\le C e^{-c b_n}\) on the circle, after
changing \(c\). Maximum modulus gives the same bound at its center.
Therefore
\[
 |f(x_n+\ii(H-d_n))|\le C m_ne^{-c b_n}\longrightarrow0
\]
by \eqref{cone:b-dominates-log}, contradicting
\eqref{cone:f-lower-center}. Both common types are impossible, proving
the existence of trivial profiles.

\medskip\noindent
\emph{A global width bound and a nontrivial limiting lift.}
Use the centers with the trivial pair in
\Cref{cor:quiet-cone}\textup{(ii)}. The bound
\eqref{cone:profile-envelope} is stronger on the coarse tail than its
hypothesis, since \(\alpha_1>1\); it may be weakened to any exponent
between \(1/2\) and \(1\). Since \(d_n/b_n=1/m_n\), we obtain
\begin{equation}\label{cone:superlinear-width}
 h_X(T/m_n)\ge H-o(d_n)
\end{equation}
uniformly for \(T\) in each fixed bounded real interval. This is a
global row width. We may therefore choose new centers \(x_n\) attaining
the row maximum that defines \(m_n\). At these centers extract limits
\(F,Q\). They are pure, and they are not both zero because
\[
 \max\{|F_n(-\ii)|,|Q_n(-\ii)|^{1/2}\}=1.
\]

Fix a real \(T\ne0\). Define, now as a function of the rescaled phase
coordinate \(Y\),
\[
 S_n(Y)=\frac{\mathfrak s(x_n+\ii H+d_nY,T/m_n)-x_n-\ii H}{d_n}.
\]
By \eqref{cone:superlinear-width}, its phase domains exhaust
\(\Im Y<0\). By the shape constraint in \Cref{thm:shape-lift},
\(\Im S_n<0\). For any \(1/2<\alpha<1\), the coarse cone estimate
anchors these half-plane-valued functions on deep compact sets:
\begin{equation}\label{cone:deep-anchor}
 |S_n(Y)-Y|\le C_TD^{1-2\alpha},\qquad D=-\Im Y\ge D_T,
\end{equation}
on the fixed compact portions under consideration, with constants
uniform in horizontal position. To obtain such uniform constants,
apply the cone theorem at each real center; the row bounds and the
lower-strip data are the same at every center.

Map the target half-plane conformally to the unit disk. Montel's theorem
then gives a subsequential limit. The anchor excludes convergence to
a boundary constant or infinity, so the limit is a holomorphic map
\(S\) of the entire lower half-plane into itself. Passing to the limit
in \eqref{cone:deep-anchor} gives the same anchor for \(S\).

\medskip\noindent
\emph{The limiting energy equation.}
Set \(h_n=T/b_n\), which tends to zero. Then
\[
 J_n(Y)=\frac{S_n(Y+h_n)-S_n(Y)}{h_n}\longrightarrow S_Y(Y)
\]
locally analytically. On a sufficiently deep half-plane, the cone
estimates and Cauchy estimates keep these steps close to \(1\) and the
fast fields small. The exact recurrence holds there with its continued
glancing root, also jointly for \(T\) along the segment from \(0\) to
the fixed value chosen. This segment, begun at \(T=0\) in the deep
small-field regime, selects the fixed sheet required by
\Cref{cone:profile-action-convergence}; no additional loop is used.
Let \(p_n(A,J)\) denote the normalized action
with rescaled shape basepoint \(A\). Its fast fields are
\[
 T F_n(A+h_nu),\qquad T^2Q_n(A+h_nu).
\]
They converge analytically to the frozen fields
\(TF(A)\), \(T^2Q(A)\). By
\Cref{cone:profile-action-convergence}, the actions and their
derivatives converge analytically to
\[
 p(A,J)=
 \begin{cases}
 p_F(TF(A),J),&Q\equiv0,\\
 p_Q(\lambda(A),J),\quad\lambda(A)^2=T^2Q(A),&F\equiv0.
 \end{cases}
\]
The second expression depends analytically on \(\lambda^2\) near zero,
so no choice of a square root is needed there.

In these coordinates \eqref{cone:normalized-recurrence} reads
\[
 p_{n,J}(S_n(Y),J_n(Y))
 -p_{n,J}(S_n(Y-h_n),J_n(Y-h_n))
 =h_n p_{n,A}(S_n(Y),J_n(Y)).
\]
Division by \(h_n\) is legitimate without any rate of convergence
assumption. Indeed, if
\(G_n(Y)=p_{n,J}(S_n(Y),J_n(Y))\), then
\[
 \frac{G_n(Y)-G_n(Y-h_n)}{h_n}
 =\int_0^1G_n'(Y-\theta h_n)\,\dd\theta,
\]
and analytic convergence gives convergence of \(G_n'\).
The limit is consequently the Euler--Lagrange equation
\[
 \partial_Yp_J(S,S_Y)=p_A(S,S_Y).
\]
For \(E(A,J)=Jp_J(A,J)-p(A,J)\), the chain rule now gives
\[
 \partial_Y E(S,S_Y)
 =S_Y\bigl(\partial_Yp_J(S,S_Y)-p_A(S,S_Y)\bigr)=0.
\]
The deep decay of the fields and the anchor, with Cauchy estimates,
give \(S_Y\to1\) as \(D\to\infty\). Hence the constant is
\begin{equation}\label{cone:limiting-energy}
 E(S,S_Y)=\frac1{12}.
\end{equation}

The stronger exponent \(\alpha_1\) now improves the anchor. At
\((\lambda,J)=(0,1)\), both frozen energies have
\(E_J=1/4\); their first field correction is quadratic by
\eqref{cone:F-series}--\eqref{cone:Q-series}. The analytic implicit
function theorem applied to \eqref{cone:limiting-energy}, together
with \eqref{cone:profile-envelope}, yields
\[
 |S_Y-1|\le C_TD^{-2\alpha_1}
\]
on a deeper half-plane. The argument of the fields is \(S(Y)\), whose
depth is comparable to \(D\) by the original anchor. Integrate along
vertical lines from infinity, where the original anchor makes
\(S-Y\) tend uniformly to zero. We obtain
\begin{equation}\label{cone:improved-anchor}
 |S(Y)-Y|\le C_TD^{1-2\alpha_1}=o(D^{-1}).
\end{equation}

\medskip\noindent
\emph{Half-plane rigidity and the contradiction.}
Put \(v(Y)=\Im(Y-S(Y))\). It is harmonic and satisfies \(v\ge-D\),
since \(S\) maps into the lower half-plane. We verify carefully that
\(v\ge0\). On a slab \(\eps<D<R\), its top boundary is bounded below
by \(-\eps\), and its bottom boundary is bounded below by
\(-\eta_R\), where
\(\eta_R=\sup_x|S(x-\ii R)-(x-\ii R)|\to0\). Compare on finite
rectangles with the harmonic function interpolating these two
constants on the horizontal sides. The extra negative contribution
needed on the vertical sides is bounded by a constant depending on
\(R\) times their harmonic measure, because \(v\ge-R\) in the slab.
Their harmonic measure tends to zero when the sides move to infinity.
Thus
\[
 v(Y)\ge-\eps\frac{R-D}{R-\eps}
              -\eta_R\frac{D-\eps}{R-\eps}.
\]
First let \(\eps\downarrow0\), then \(R\to\infty\), to get
\(v\ge0\). Harnack's inequality in the half-plane gives, for
\(R\ge1\),
\[
 v(-\ii R)\ge R^{-1}v(-\ii).
\]
For example, this is the ordinary disk Harnack inequality after
mapping \(-\ii\) to the origin, since \(-\ii R\) has disk modulus
\((R-1)/(R+1)\). By \eqref{cone:improved-anchor} its left side is
\(o(R^{-1})\). Hence \(v(-\ii)=0\), and the strong minimum principle
gives \(v\equiv0\). The holomorphic function \(S-Y\) is therefore real
constant; the anchor makes that constant zero. We have proved
\(S(Y)=Y\).

Finally, set \(J=1\) in \eqref{cone:limiting-energy}. Deep down the
field parameter is small, and the energy expansions are
\[
 E_F(\lambda,1)-\frac1{12}
 =-\frac{\lambda^2}{480}+O(\lambda^4),\qquad
 E_Q(\lambda,1)-\frac1{12}
 =-\frac{\lambda^2}{60}+O(\lambda^4).
\]
Each right side can vanish near zero only when \(\lambda^2=0\).
Since \(T\ne0\), the surviving pure field vanishes on a deep open
set, and hence identically. This contradicts its normalization at
\(-\ii\), proving the Theorem.
\end{proof}
